% ==============================================================================
% 249_FFGFT_Falsifizierbarkeit_De_ch.tex
% Kapitelinhalt – einzubinden via Wrapper
% Autor: Johann Pascher · Mai 2026
% ==============================================================================

\section*{Abstract}

Dieses Dokument ist ein Begleitdokument zu Dok.~248 (FFGFT --- Epistemologische Position). Es klärt die wissenschaftstheoretische Position von FFGFT gegenüber metaphysischen Spekulationen: FFGFT ist eine physikalische Theorie mit empirischem Gehalt und konkreten Falsifizierungsbedingungen. Das Postulat einer ``reinen, ungebundenen Information'' als Primitiv (Myser 2026) ist dagegen eine metaphysische Rahmung ohne Ableitungsmacht --- nicht widerlegbar, nicht anschlussfähig an bekannte Physik. Der entscheidende Unterschied liegt nicht nur in der Wahl des Primitivs, sondern in der Konsequenz aus dieser Wahl.

\tableofcontents
\newpage

% ------------------------------------------------------------------------------
\section{Falsifizierbarkeit als wissenschaftliches Kriterium}

Eine wissenschaftliche Theorie unterscheidet sich von einer metaphysischen Spekulation durch ihre \textbf{Anfälligkeit für Widerlegung} \cite{Popper1934}. Eine Theorie, die durch kein denkbares Experiment falsifiziert werden kann, macht keine physikalischen Aussagen --- sie macht Aussagen über das, was wir für möglich halten, nicht über das, was die Welt ist.

Dieses Kriterium ist kein Dogma, sondern ein methodologisches Werkzeug: Es zieht eine klare Linie zwischen Theorien, die sich der Welt aussetzen, und solchen, die sich ihr entziehen.

\begin{keyresult}[Poppers Kriterium]
Eine Theorie ist wissenschaftlich, wenn sie \textbf{prinzipiell falsifizierbar} ist: wenn es Beobachtungen oder Experimente gibt, die \emph{im Widerspruch} zur Theorie stünden und sie damit widerlegen würden. Eine Theorie, die mit jeder möglichen Beobachtung kompatibel ist, sagt über die Welt nichts aus.
\end{keyresult}

% ------------------------------------------------------------------------------
\section{FFGFT ist falsifizierbar}

FFGFT macht konkrete, prädiktive Aussagen, die prinzipiell widerlegt werden könnten.

\subsection{Minimale Längenskala $L_0$}

\begin{equation}
  L_0 = \xipar \cdot \lP \approx 2{,}2 \times 10^{-39} \text{ m}
\end{equation}

Diese absolute Untergrenze der Raumzeit-Struktur ist eine Vorhersage von FFGFT. Würden zukünftige Hochenergie-Experimente oder kosmologische Beobachtungen konsistent Strukturen \emph{unterhalb} dieser Skala zeigen, wäre FFGFT widerlegt. Umgekehrt: Blieben Strukturen oberhalb dieser Grenze \emph{immer} Planck-skalig oder größer, stützt das die Vorhersage.

\subsection{Fraktale Dimension $D_f$}

Die fraktale Dimension der Raumzeit in FFGFT beträgt:
\begin{equation}
  D_f = 2{,}999867
\end{equation}
Das ist eine Abweichung von exakt 3 um den Faktor $1{,}33 \times 10^{-4}$. Diese Abweichung ist prinzipiell messbar über:
\begin{itemize}
  \item \textbf{Gammablitze-Laufzeitdispersion}: Energieabhängige Lichtlaufzeit aus kosmologischen Quellen ist ein etabliertes Suchfeld für Quantengravitationseffekte. FFGFT sagt eine spezifische, $\xipar$-abhängige Dispersionsstruktur voraus.
  \item \textbf{Gravitationswellen-Phasenstruktur}: Abweichungen in der Wellendispersion bei extrem hohen Frequenzen.
\end{itemize}
Würde eine solche Messung eine Abweichung zeigen, die nicht mit $D_f = 2{,}999867$ kompatibel ist, wäre FFGFT falsifiziert.

\subsection{Keine Singularität in schwarzen Löchern}

FFGFT sagt vorher, dass schwarze Löcher keine Singularität enthalten --- weil $L_0$ eine absolute untere Dichtegrenze setzt, die strukturell nicht unterschritten werden kann.

Diese Vorhersage ist unterscheidbar von der allgemeinen Relativitätstheorie durch:
\begin{itemize}
  \item \textbf{Gravitationswellen-Echos}: Quantengravitationsmodelle ohne Singularität sagen charakteristische Echosignale nach Verschmelzungen voraus, die LIGO/Virgo in zukünftigen Generationen nachweisen könnte.
  \item \textbf{Hawking-Strahlungsspektrum}: FFGFT sagt ein modifiziertes Spektrum voraus, das die $T^4$-Windungsstruktur reflektiert --- unterscheidbar vom rein thermischen Spektrum der Standardtheorie.
\end{itemize}

\subsection{Kopplungskonstanten aus $\xipar$}

FFGFT leitet alle fundamentalen Kopplungskonstanten aus einem einzigen Parameter $\xipar = 1{,}333 \times 10^{-4}$ ab. Jede experimentelle Verfeinerung dieser Konstanten ist ein Test von FFGFT:
\begin{itemize}
  \item Feinstrukturkonstante $\alpha^{-1} = 137{,}036$ --- FFGFT-Ableitung: $\alpha^{-1} \approx 137{,}036$ (Abweichung $< 0{,}001\%$)
  \item Anomales magnetisches Moment des Elektrons $a_e$ --- FFGFT-Korrektur: $K_\text{frak}^{3/2}$ reduziert Abweichung von $2{,}03\%$ auf $0{,}014\%$
  \item Koide-Formel für Lepton-Massen --- FFGFT liefert geometrische Ableitung
\end{itemize}
Jede dieser Ableitungen kann durch präzisere Messungen bestätigt oder widerlegt werden.

% ------------------------------------------------------------------------------
\section{Das Postulat ``reiner Information'' ist nicht falsifizierbar}

Das Postulat von Atticus Myser --- ``pure unrestrained information'' als Primitiv, ohne Struktur, ohne Unterscheidung, ohne Bindung --- macht \textbf{keine prädiktiven Aussagen}.

\subsection{Keine ableitbaren Größen}

Aus ``pure unrestrained information'' folgt keine Messgröße, keine Längenskala, keine Kopplungskonstante, keine Feldgleichung. Das Postulat beschreibt, was \emph{vor} aller Struktur vermutet wird --- aber ohne Struktur gibt es keine Ableitung, und ohne Ableitung gibt es keine Vorhersage.

\subsection{Kompatibel mit jeder Beobachtung}

Da aus dem Postulat nichts folgt, kann es auch durch nichts widerlegt werden. Jede Beobachtung --- ob sie für FFGFT, für Stringtheorie, für MOND, oder für die Standardkosmologie spricht --- ist mit ``pure unrestrained information'' kompatibel, weil das Postulat keine Beobachtung ausschließt.

\begin{wichtig}[Nicht falsifizierbar = nicht physikalisch]
Eine Theorie, die mit jeder Beobachtung kompatibel ist, macht keine physikalischen Aussagen. Das Postulat der ``reinen Information'' ist daher keine physikalische Theorie --- es ist eine \textbf{metaphysische Rahmung}. Das ist nicht per se wertlos, aber es ist ein grundlegend anderer Anspruch als der von FFGFT.
\end{wichtig}

\subsection{Kein Occlusion-Mechanismus ableitbar}

Atticus' eigene Theorie der ``O-bits'' und der sequenziellen Adhäsion (Occlusion) macht durchaus interessante Strukturaussagen. Aber diese folgen nicht aus dem Postulat der ``reinen Information'' --- sie werden \emph{zusätzlich} eingeführt, als separate Regel. Diese Regel ist selbst bereits Struktur. Das Postulat des Primitivs trägt die Last der Occlusion nicht --- sie wird von außen eingetragen.

% ------------------------------------------------------------------------------
\section{Der wissenschaftstheoretische Unterschied}

FFGFT und Atticus' Ansatz unterscheiden sich nicht nur in der Wahl des Primitivs ($T^4$ vs.\ ``pure information''), sondern in einem tieferen methodologischen Punkt:

\begin{longtable}{>{\bfseries}p{3.5cm} p{5cm} p{5cm}}
\toprule
Kriterium & FFGFT & ``Pure information'' (Myser) \\
\midrule
Primitiv & $T^4 + \xipar$ (geometrisch, posiert) & ``pure unrestrained information'' (posiert) \\
Operational & Ja --- Derivationen, Feldgleichungen, Messung & Nein --- keine Ableitung möglich \\
Falsifizierbar & Ja --- $L_0$, $D_f$, Singularitäten, $\alpha$ & Nein --- keine ausgeschlossene Beobachtung \\
Vorhersagen & Konkret, numerisch & Keine \\
Anschlussfähig & An Standardmodell, GR, QFT & Nicht direkt \\
Wissenschaftstyp & Physikalische Theorie & Metaphysische Rahmung \\
\bottomrule
\end{longtable}

\begin{erkenntnis}[Der entscheidende Unterschied]
\textbf{FFGFT ist eine physikalische Theorie mit empirischem Gehalt.} Das Postulat der ``reinen Information'' ist eine metaphysische Spekulation ohne Ableitungsmacht. Der entscheidende Unterschied liegt nicht nur in der Axiomensetzung, sondern in der \emph{Konsequenz} aus ihr: Was folgt aus dem Primitiv? Wie viel folgt? Und kann das, was folgt, durch Beobachtung widerlegt werden?
\end{erkenntnis}

% ------------------------------------------------------------------------------
\section{Falsifizierbarkeit und epistemische Offenheit}

Diese Unterscheidung ist kein Angriff auf metaphysische Rahmungen. Metaphysische Fragen --- warum gibt es überhaupt etwas? was ist das Wesen von Information? was liegt vor aller Struktur? --- sind legitime und wichtige Fragen. Sie sind nur nicht durch physikalische Experimente entscheidbar.

FFGFT stellt keine Antwort auf diese Fragen bereit. Es lässt ``warum $T^4$?'' explizit offen (vgl.\ Dok.~248, Abschnitt 8). Das ist strukturelle Ehrlichkeit: FFGFT beschreibt \emph{von innen}, mit den Mitteln, die eine Innenbeschreibung zur Verfügung hat.

Was FFGFT von metaphysischen Spekulationen unterscheidet, ist nicht das Fehlen einer Antwort auf die ultimative Frage --- sondern das Vorhandensein konkreter, widerlegbarer Aussagen über die Welt, die wir sehen.

\begin{erkenntnis}
\textbf{FFGFT ist falsifizierbar, aber epistemisch bescheiden.} Es macht keine Aussagen über das, was vor $T^4$ liegt. Es macht präzise Aussagen über das, was aus $T^4 + \xipar$ folgt --- und diese Aussagen können durch Beobachtung widerlegt werden. Das ist der Anspruch einer physikalischen Theorie, nicht mehr und nicht weniger.
\end{erkenntnis}

% ------------------------------------------------------------------------------
\section*{Zusammenfassung}

\begin{longtable}{>{\bfseries}p{7.2cm} p{7.2cm}}
\toprule
Aussage & Position \\
\midrule
FFGFT macht keine prädiktiven Aussagen & Falsch --- $L_0$, $D_f$, Singularitäten, $\alpha$ sind Vorhersagen \\
FFGFT ist nicht falsifizierbar & Falsch --- jede der Vorhersagen ist prinzipiell widerlegbar \\
``Pure information'' ist eine physikalische Theorie & Nein --- metaphysische Rahmung ohne Ableitungsmacht \\
``Pure information'' ist falsifizierbar & Nein --- keine Beobachtung ist ausgeschlossen \\
Beide Ansätze sind wissenschaftstheoretisch gleichwertig & Nein --- nur FFGFT ist eine physikalische Theorie im Popperschen Sinne \\
Metaphysische Rahmungen sind wertlos & Nein --- sie sind legitim, aber von anderer Art als physikalische Theorien \\
FFGFT beantwortet ``warum $T^4$?'' & Nein --- diese Frage bleibt explizit offen (vgl.\ Dok.~248) \\
\bottomrule
\end{longtable}
